% =====================================================================
%  CARNET D'ENTRAINEMENT — FICHE 6 — THÈME 4
%  Niveau FACILE — Exercices
% =====================================================================
\documentclass[11pt,a4paper]{article}
\usepackage[utf8]{inputenc}
\usepackage[T1]{fontenc}
\usepackage{lmodern}
\usepackage[french]{babel}
\usepackage{amsmath,amssymb,mathtools}
\usepackage{icomma}
\usepackage[version=4]{mhchem}
\usepackage{siunitx}
\sisetup{output-decimal-marker={,}}
\usepackage{xcolor}
\usepackage{colortbl}
\usepackage{tikz}
\usepackage[most]{tcolorbox}
\usepackage{enumitem}
\usepackage{array}
\usepackage{booktabs}
\usepackage{fancyhdr}
\usepackage[hidelinks]{hyperref}
\usetikzlibrary{arrows.meta,positioning,calc,shapes.geometric}
\usepackage[a4paper,margin=2.1cm,headheight=26pt,headsep=8pt]{geometry}

\definecolor{primary}{RGB}{146,95,160}
\definecolor{primarydark}{RGB}{92,52,104}
\definecolor{excol}{RGB}{0,135,90}
\definecolor{atomC}{RGB}{80,80,88}
\definecolor{atomH}{RGB}{235,235,235}
\definecolor{atomO}{RGB}{220,55,45}
\definecolor{atomN}{RGB}{48,80,200}
\definecolor{lonepair}{RGB}{120,94,200}
\definecolor{bondcol}{RGB}{60,60,70}

\newcommand{\ball}[5]{\shade[ball color=#2] (#1) circle (#3);\node[font=\footnotesize\bfseries,text=#4] at (#1){#5};}
\newcommand{\pbond}[2]{\draw[bondcol,line width=1.6pt] (#1)--(#2);}
\newcommand{\wbond}[2]{\draw[bondcol,fill=bondcol,draw=none] let \p1=($(#2)-(#1)$), \n1={atan2(\y1,\x1)} in (#1) -- ($(#2)+(\n1+90:0.12)$) -- ($(#2)+(\n1-90:0.12)$) -- cycle;}
\newcommand{\hbond}[2]{\draw[bondcol,line width=1pt,dash pattern=on 1.4pt off 1.8pt] (#1)--(#2);}
\newcommand{\lobe}[3]{\begin{scope}[shift={(#1)},rotate=#2]\shade[ball color=lonepair!55,opacity=0.9] (#3,0) ellipse (0.34 and 0.20);\end{scope}}
\newcommand{\AX}[2]{\ensuremath{\mathrm{AX_{#1}E_{#2}}}}

\tcbset{boxbase/.style={enhanced,breakable,boxrule=0.9pt,arc=2.5pt,
   left=3mm,right=3mm,top=2.2mm,bottom=2.2mm,
   fonttitle=\bfseries,coltitle=white,toptitle=0.6mm,bottomtitle=0.6mm}}
\newtcolorbox[auto counter]{exo}[1][]{boxbase,colback=primary!5,colframe=primary,
   title={\textsc{Exercice}~\thetcbcounter\ifx\relax#1\relax\relax\else\textmd{ —\itshape\ #1}\fi}}

\pagestyle{fancy}\fancyhf{}
\fancyhead[L]{\small\textcolor{primarydark}{\textbf{Fiche 6 · Thème 4} — VSEPR et géométrie moléculaire}}
\fancyhead[R]{\small\textcolor{primarydark}{\textsc{Facile}}}
\fancyfoot[C]{\small\thepage}
\renewcommand{\headrulewidth}{0.8pt}
\renewcommand{\headrule}{\hbox to\headwidth{\color{primary}\leaders\hrule height \headrulewidth\hfill}}
\setlength{\parindent}{0pt}\setlength{\parskip}{3pt}

\begin{document}
\begin{center}
{\Large\bfseries\color{primarydark}Thème 4 — Niveau Facile}\\[2pt]
{\small\itshape Lire la fiche \texttt{Tuto\_Theme-4} avant de commencer. Aucune calculatrice n'est nécessaire.}
\end{center}
\medskip

\begin{exo}[écrire \AX{n}{m} et la géométrie de répulsion]
On donne, pour un atome central A, le nombre $n$ d'atomes liés et le nombre $m$ de doublets non
liants. Pour chaque cas : écrire le symbole \AX{n}{m}, donner le \textbf{total de sites} $n+m$ et
en déduire la \textbf{géométrie de répulsion}.
\begin{enumerate}[label=\alph*),leftmargin=2em,itemsep=1pt]
  \item $n=2$, $m=0$ ;
  \item $n=3$, $m=0$ ;
  \item $n=4$, $m=0$ ;
  \item $n=2$, $m=2$.
\end{enumerate}
\end{exo}

\begin{exo}[lire la géométrie moléculaire dans la table]
À l'aide du tableau de synthèse VSEPR, donner la \textbf{géométrie moléculaire} correspondant à
chacun des types suivants :
\[ \AX{2}{0}, \qquad \AX{3}{0}, \qquad \AX{4}{0}, \qquad \AX{3}{1}, \qquad \AX{2}{2}. \]
\end{exo}

\begin{exo}[compter \texorpdfstring{$n$}{n} et \texorpdfstring{$m$}{m} autour de l'atome central]
Pour chaque molécule décrite, compter le nombre $n$ d'atomes liés à l'atome central et le nombre
$m$ de doublets non liants qu'il porte.
\begin{enumerate}[label=\alph*),leftmargin=2em,itemsep=2pt]
  \item \ce{CH4} : le carbone forme 4 liaisons simples \ce{C-H} et n'a aucun doublet libre ;
  \item \ce{NH3} : l'azote forme 3 liaisons \ce{N-H} et porte 1 doublet libre ;
  \item \ce{H2O} : l'oxygène forme 2 liaisons \ce{O-H} et porte 2 doublets libres ;
  \item \ce{CO2} (\ce{O=C=O}) : le carbone forme 2 doubles liaisons et n'a aucun doublet libre.
\end{enumerate}
\end{exo}

\begin{exo}[associer chaque géométrie à son angle idéal]
Relier chaque géométrie de la colonne de gauche à l'angle idéal de la colonne de droite.
\begin{center}
\renewcommand{\arraystretch}{1.3}
\begin{tabular}{r@{\qquad}l}
\textbf{Géométrie} & \textbf{Angle idéal}\\
\hline
1. linéaire            & A. $109{,}5^\circ$\\
2. triangulaire plane  & B. $180^\circ$\\
3. tétraédrique        & C. $120^\circ$\\
\end{tabular}
\end{center}
\end{exo}

\begin{exo}[vrai ou faux ?]
Indiquer si chaque affirmation est \textbf{vraie} ou \textbf{fausse}, et justifier en une phrase.
\begin{enumerate}[label=\alph*),leftmargin=2em,itemsep=2pt]
  \item « Une double liaison compte pour deux sites de répulsion. »
  \item « Quand l'atome central ne porte aucun doublet libre, la géométrie de répulsion est
        identique à la géométrie moléculaire. »
\end{enumerate}
\end{exo}

\end{document}
